% ---- pg-479 （书页 467） ----

\[
\mathbf{X}_1=\begin{bmatrix}2&12\\4&10\\3&8\end{bmatrix};\qquad \bar{\mathbf{x}}_1=\begin{bmatrix}3\\10\end{bmatrix},\quad 2\mathbf{S}_1=\begin{bmatrix}2&-2\\-2&8\end{bmatrix}
\]
\[
\mathbf{X}_2=\begin{bmatrix}5&7\\3&9\\4&5\end{bmatrix};\qquad \bar{\mathbf{x}}_2=\begin{bmatrix}4\\7\end{bmatrix},\quad 2\mathbf{S}_2=\begin{bmatrix}2&-2\\-2&8\end{bmatrix}
\]

联合协方差阵为
\[
\mathbf{S}_{\mathrm{p}}=\frac{1}{4}(2\mathbf{S}_1+2\mathbf{S}_2)=\begin{bmatrix}1&-1\\-1&4\end{bmatrix}
\]
利用 $\mathbf{S}_{\mathrm{p}}$, 其余数据以及\hyperref[eq:11-18]{式(11-18)}中同等代价同等先验概率的法则, 我们就可以对样本观测值进行分类. 然后可以验证(见练习 11.19), 模糊矩阵为
\[
\begin{array}{r|cc}
 & \multicolumn{2}{c}{\text{分类为:}}\\ \hline
\text{真实总体:} & \pi_1 & \pi_2\\ \hline
\pi_1 & 2 & 1\\
\pi_2 & 1 & 2
\end{array}
\]
因此,
\[
\mathrm{APER}(\text{表观失误率})=2/6=0.33
\]

从 $\mathbf{X}_1$ 中提留第一个观测值 $\mathbf{x}_H'=[2,12]$, 计算
\[
\mathbf{X}_{1H}=\begin{bmatrix}4&10\\3&8\end{bmatrix};\qquad \bar{\mathbf{x}}_{1H}=\begin{bmatrix}3.5\\9\end{bmatrix};\qquad 1\mathbf{S}_{1H}=\begin{bmatrix}0.5&1\\1&2\end{bmatrix}
\]
新的联合协方差阵 $\mathbf{S}_{H,\mathrm{p}}$ 为
\[
\mathbf{S}_{H,\mathrm{p}}=\frac{1}{3}[1\mathbf{S}_{1H}+2\mathbf{S}_2]=\frac{1}{3}\begin{bmatrix}2.5&-1\\-1&10\end{bmatrix}
\]
其逆\footnote{巴特利特(见文献[3])给出一个矩阵恒等式, 可用来直接从 $\mathbf{S}_{\mathrm{p}}^{-1}$ 迅速算出 $\mathbf{S}_{H,\mathrm{p}}^{-1}$. 这样就不需要在取出每个观测值后重新计算逆矩阵(见练习 11.20).}为
\[
\mathbf{S}_{H,\mathrm{p}}^{-1}=\frac{1}{8}\begin{bmatrix}10&1\\1&2.5\end{bmatrix}
\]
根据组均值 $\bar{\mathbf{x}}_{1H}$ 和 $\bar{\mathbf{x}}_2$ 之间的样本距离来分类提留观测值 $\bar{\mathbf{x}}_{1H}$, 计算起来更快一些. 这种方法等价于计算线性函数 $\hat{y}=\hat{\mathbf{a}}_{H}'\mathbf{x}_H=(\bar{\mathbf{x}}_{1H}-\bar{\mathbf{x}}_2)'\mathbf{S}_{H,\mathrm{p}}^{-1}\mathbf{x}_H$ 的值, 然后将它与中点 $\hat{m}_{H}=\frac{1}{2}(\bar{\mathbf{x}}_{1H}-\bar{\mathbf{x}}_2)'\mathbf{S}_{H,\mathrm{p}}^{-1}(\bar{\mathbf{x}}_{1H}+\bar{\mathbf{x}}_2)$ 作比较[见\hyperref[eq:11-19]{式(11-19)}和\hyperref[eq:11-20]{式(11-20)}].

这样, 对 $\mathbf{x}_H'=[2,12]$ 我们有

到 $\bar{\mathbf{x}}_{1H}$ 的平方距离 $=(\mathbf{x}_H-\bar{\mathbf{x}}_{1H})'\mathbf{S}_{H,\mathrm{p}}^{-1}(\mathbf{x}_H-\bar{\mathbf{x}}_{1H})$
\[
=[2-3.5\quad 12-9]\frac{1}{8}\begin{bmatrix}10&1\\1&2.5\end{bmatrix}\begin{bmatrix}2-3.5\\12-9\end{bmatrix}=4.5
\]
到 $\bar{\mathbf{x}}_2$ 的平方距离 $=(\mathbf{x}_H-\bar{\mathbf{x}}_2)'\mathbf{S}_{H,\mathrm{p}}^{-1}(\mathbf{x}_H-\bar{\mathbf{x}}_2)$
